% !TeX root = ../main.tex
\section{Group Cohomology of Surface Groups}
\label{sec:L04-W02D2}

\begin{example}[Compactly supported cohomology of \(R^n\)]
    \label{ex:compact-supported-cohomology-Rn}
    Suppose \(\R^n\) is equipped with a polyhedral decomposition. 
    Suppose each \((n-1)\)-cell is adjacent to \(2\) oriented \(n\)-cells, and the \(n-1\)-cell appears in their boundary with opposite orientation. 
    Suppose further that each compact \(K\subseteq \R^n\) is contained in a ball. (For example, regard \(R^3\) as  the universal covering space of 3-\(torus\).)

    Since \(\R^n\) is contractible, \(H_c^0(\R^n;\Z)=0\).

    For \(1\leq i <n\), take \(\eta \in C_c^i(\R^n;\Z)\) such that \(\d\eta=0\). Take a ball \(B\) such that \( \supp{\eta} \subseteq B \). We can regard \(\eta\) as an element of \(H^i(\R^n,\R^n\backslash B)\) (that is, a cochain on \(\R^n\) that vanishes on \(B\)). Recall that \(H^i(\R^n,\R^n\backslash B) \cong \tilde H^{i-1}(\R^n\backslash B) \cong \tilde H^{i-1}(\mathbb{S}^{n-1}) =0 \).

    Consider now the augmented chain \(C_c^{n-1}(\R^n;\Z) \xrightarrow{\d} C_c^n(\R^n;\Z) \xrightarrow{\epsilon} \Z \to 0\). The assumption that each \((n-1)\)-cell occurs in the boundary of two \(n\)-cells with opposite orientation implies that \(\im \d \subseteq \ker \epsilon\) (similar to \cref{ex:compact-supported-cohomology-R}), so \(\epsilon\) factors through \(H_c^n(\R^n;\Z) \to \Z\). Again, we need to show that the map is injective. The argument is similar to \cref{ex:compact-supported-cohomology-R}: we take the dual graph of \(\R^n\) so that vertices are dual \(n\)-cells and edges are dual \((n-1)\)-cells. We then adjust one dual \(n\)-cell at a time until the cochain is supported on a single dual \(n\)-cell with coefficient zero.
\end{example}

\begin{corollary}
    \label{cor:cohomology-of-surface-group}
    Let \(M\) be a closed \(n\)-manifold with \(\tilde M \cong \R^n\). Let \(G=\pi_1(M)\). Then 
    \begin{align*}
        H^i(G;\Z G) \cong H^i_c(\R^n;\Z) = \begin{cases}
            \Z & i=n,\\
            0 & \text{otherwise}.
        \end{cases}
    \end{align*}

    In particular, if \(\Sigma_g\) is a closed orientable surface with genus \(g \geq 1\) and \(G=\pi_1(\Sigma_g)\), then 
    \begin{align*}
        H^i(G;\Z G) = \begin{cases}
            \Z & i=2,\\
            0 & \text{otherwise}.
        \end{cases}
    \end{align*}
\end{corollary}
\begin{remark}
    For non-orientable manifolds, the cohomology would be the same. What's different is the action of \(\G\) on the top cohomology \(Z\).
\end{remark}

In general, \(H^i(G;\Z G)\) is a \(\Z G\)-module. Take a free resolution \(F\). A priori, \(\Hom{F_i }{\Z G }{\Z G}\) is just an abelian group, but the right action by \(G\) satisfies  the criterion from \cref{lem:compact-homomorphism-cong-ZG-homomorphism}. Therefore, as a quotient of \(\Z G\)-module by \(\Z G\)-module, \(H^i(G;\Z G)\) is also a \(\Z G\)-module.

\begin{question}[Wall]
    \label{q:wall}
    Suppose \(G\) has type \(F\) (recall \cref{def:finite-type-KG1}) and suppose that for some \(n\), \(H^n(G;\Z G) = \Z\) and \(H^i(G; \Z G) =0\) otherwise. Does there exist a closed aspherical manifold \(M\) such that \(G=\pi_1 (M)\)?
\end{question}

\cref{q:wall} is known to be true for \(n= 1\) (\(G =\Z\))  and \(n=2\) (surface groups). It is wide open for \(n\geq 3\). 

Notice that if such an \(M\) exists, it's unique up to homotopic equivalence by the uniqueness of \(K(G,1)\). It's natural to ask:

\begin{question}[Borel conjecture]
    \label{q:borel-conjecture}
    If \(M_1, M_2\) are closed aspherical manifolds that are homotopic equivalent to each other, are they homeomorphic to each other? 
\end{question}

\cref{q:borel-conjecture} is known to be true for \(n=1,2,3\). Note that in general, without the assumption of aspherical, this is not true.

Recall that two homotopic equivalent manifolds \(M_1, M_2\) have isomorphic fundamental groups. 
A group theoretic version of \cref{q:borel-conjecture} is: Given a group \(G\), does it have a unique \(K(G,1)\) up to homeomorphism?



\begin{theorem}
    \label{thm:n-connected-gives-aspherical-cw-complex}
    Let \(X\) be a CW complex such that \(\tilde X\) is \(k\)-connected (recall \cref{def:n-connected}) for some \(k\). Then there exists an aspherical CW complex \(Z\) such that \(Z^{k+1} = X^{k+1}\).
\end{theorem}

\cref{thm:n-connected-gives-aspherical-cw-complex} is the reason we are interested in \(1\)-connected spaces, because we can obtain a CW complex with the same fundamental group.

\begin{proof}
    Assume inductively that we've built \(Z_n\) such that \(\tilde Z_n\) is \(n\)-connected and \(Z_n^{k+1} = X^{k+1}\) (\(Z_k\) is just \(X\)). Note that every map \(f:\mathbb{S}^{n+1}\to Z_n\) is homotopic to a map with image in \(Z_n^{n+1}\). Let \(\{F_\a : \mathbb{S}^{n+1} \to Z_n^{n+1} \mid  \a \in A_n\}\) denote the homotopy classes of maps \(\mathbb{S}^{n+1}\to Z_n\). Build \(Z_{n+1}\) by gluing \((n+1)\)-cells to \(Z_n\) using  \(F_\a\)s as attaching maps.

    Observe that \(Z_{n+1}^{n+1} = Z_n^{n+1}\). Any \(g: \mathbb{S}^{n+1} \to Z_{n+1}\) is homotopic to some \(F_\a\), hence is null-homotopic. Therefore \(\tilde Z_{n+1}\) is \((n+1)\)-connected.
\end{proof}